\scanpage{72}
% archive foliation: 71
% lines 1--4
 de $X_i$ suivant \unsure{ses|les} points à l'infini\ldots).
Cl. que $\widehat X_i$ est automatiquement $\omega$-orienté,
donc aussi
\[\widehat X=\coprod\widehat X_i\]
% lines 5--12
De cette façon, la catégorie envisagée
(des surfaces topologiques ``spéciales''
$\omega$-\unsure{orientation|orientées}\inserted{-- aux homéomorphismes $\omega$-orientés}) devient équivalente à
celle des surfaces topologiques \inserted{$\omega$-orientées}
à comp. conn. compactes \uline{munies}
d'une partie \deleted{\unsure{fi}} $S$ discrète,
avec les isom. de couples $(\widehat X,S)$.
% lines 13--20
De $\overline{\mathrm m}$, la catégorie des couples
\deleted{$\unsure{(X,\widehat X)|(X,o)}$} $(X,o,J)$ avec $(X,o)$ comme
dessus ($o$ la $\omega$-orientation) et
$J$ partie discrète de $X$, est équivalente à celle des systèmes
$(\widehat X,o,S,J)$, où $(\widehat X,o)$ est une surface
$\omega$-orientée à comp. conn. compactes,
